% !PRINTED 384
\weberpage{384}

%% [running head, not body text -- commented out 2026-08-07] Page furniture (running
%% head, folio, signature) is excluded from this edition by standing policy.
%% \noindent 384 \hfill \textsc{Elfter Abschnitt.} \hfill \S.\,124.

%<<CONT
% [source-reviewed 2026-08-26] The faint mark before the first word at the top edge is a scan speck, not body text; no text is transcribed.
\gesperrt{$\omega$ ist positiv und grösser als 1, und die mit $\omega$ conjugirte Zahl $\omega'$
ist negativ und dem absoluten Werth nach kleiner als 1,}

oder in Zeichen:

\weq{1}{0 < y\sqrt{d} - x < 1 < y\sqrt{d} + x.}

Deutet man in einer Ebene $x, y$ als rechtwinklige Coordinaten, so
dass die Punkte der Ebene die Bilder der Zahlen $\omega$ werden, so liegen
also die Bilder der reducirten Zahlen in dem Theile der Ebene, der in
der Fig. 27, die der Annahme $d = 2$ entspricht, schraffirt ist.

% SOURCE FIGURE CROP: Bd. I registered scan 410 (printed p. 384), 600 dpi,
% crop x=440,y=1230,w=1160,h=870; source-faithful raster, Fig. 27.
\weberfigureimage{weber-b1-fig27.png}{0.48}

Wir nehmen hier eigentliche und uneigentliche Aequivalenz zusammen
und beweisen zunächst die Sätze:

\noindent 1.\quad \gesperrt{Jede Zahl $\omega$ ist mit einer reducirten Zahl äquivalent.}

Man sieht dies leicht ein, wenn man $\omega$ in einen Kettenbruch entwickelt
und bis zur Schlusszahl $\omega_n$ geht:

\weq{2}{\omega = (\alpha_0, \alpha_1, \alpha_2 \ldots \alpha_{n-1}, \omega_n).}

Dann ist $\omega$ mit $\omega_n$ äquivalent, und $\omega_n$ ist, sobald $n$ grösser als 0
ist, positiv und grösser als 1. Da $\omega_n$ mit $\omega$ äquivalent ist, so handelt
es sich nur noch um den Nachweis, dass, wenn $n$ hinlänglich gross ist,
die zu $\omega_n$ conjugirte Zahl $\omega'_n$ negativ und absolut kleiner als 1 ist.
Dies aber ersieht man aus nachstehenden Formeln: Aus (2) folgt, wenn
$P_n : Q_n$ die Näherungsbrüche sind,
\[
\omega = \dfrac{P_n \omega_n + P_{n-1}}{Q_n \omega_n + Q_{n-1}},
\]
oder durch Auflösung

\weq{3}{\omega_n = -\dfrac{Q_{n-1}\omega - P_{n-1}}{Q_n\omega - P_n} = -\dfrac{Q_{n-1}}{Q_n} - \dfrac{(-1)^n}{Q_n(Q_n\omega - P_n)}.}

Verwandeln wir in dieser Gleichung $\sqrt{d}$ in $-\sqrt{d}$, so geht gleichzeitig
$\omega$ in $\omega'$, $\omega_n$ in $\omega'_n$ über und wir erhalten aus (3)

\weq{4}{\omega'_n = -\dfrac{Q_{n-1}\omega' - P_{n-1}}{Q_n\omega' - P_n} = -\dfrac{Q_{n-1}}{Q_n} \dfrac{\omega' - \dfrac{P_{n-1}}{Q_{n-1}}}{\omega' - \dfrac{P_n}{Q_n}},}
%CONT>>
